\chapter*{Útmutató és jelölések}
\addcontentsline{toc}{chapter}{Útmutató és jelölések}
A jegyzet először a sima rendszerek stabilitását ismétli át, majd a kapcsolás, a sliding, az ütközés és a grazing hatását tárgyalja. Végül a periodikus pályák numerikus számítása és követése következik. Általános háttér: \cite{dib}.

A folytonos idő szerinti deriváltat pont, a normált szegmensidő szerinti deriváltat vessző jelöli. A vektorok oszlopvektorok; a skalár függvény deriváltja sorvektor, $H_x=(\nabla H)^\mathsf T$. A transzponálást $\mathsf T$ jelöli.
\begin{longtable}{@{}p{.22\textwidth}p{.72\textwidth}@{}}
\toprule Jelölés & Jelentés\\\midrule\endhead
$x,z$ & Állapotvektor; a mechanikai példákban pozíciókat és fizikai sebességeket tartalmaz.\\
$F,f$ & Folytonos vektormező; a diszkrét modelleknél $f$ leképezést is jelöl, amit a szöveg külön megad.\\
$\Phi(x,t)$ & A folytonos rendszer megoldásoperátora.\\
$X(t),N_T$ & Fundamentális, illetve időbeli állapotátmeneti mátrix.\\
$\mathcal M,\mathbf M$ & Egy periodikus pálya monodrómia-mátrixa.\\
$\lambda$ & Az adott szövegkörnyezetben Jacobi-sajátérték vagy Floquet-multiplikátor; a stabilitási feltétel ezekre eltérő.\\
$H,\Sigma$ & Eseményfüggvény és kapcsolófelület: $\Sigma=\{H=0\}$.\\
$H_x,H_{xx}$ & Az eseményfüggvény sorvektor deriváltja és Hess-mátrixa.\\
$v_H,a_H$ & $H$ első és második időderiváltja a sima pálya mentén.\\
$R,Q_x$ & Reset leképezés és a közös időre vonatkozó korrekciós mátrix.\\
$\Pi,P$ & Poincaré-metszet és visszatérési leképezés.\\
$\sigma,r,\zeta$ & Falhelyzet, ütközési tényező és dimenziótlan viszkózus csillapítás.\\
$T,T_1,T_2,Q$ & Teljes periódus, szegmensidők és a harmonikus példa negyed periódusa.\\
$\mu,\Delta s$ & Fizikai vagy lokális modellparaméter és pseudo-arclength lépéshossz.\\
$\nu,\chi$ & A skaláris square-root modell kontrakciós paramétere és a lokális távolságeltolás.\\
$\hat E,\kappa$ & Fajlagos mechanikai energia és fajlagos lábmerevség: $[\hat E]=\mathrm{J/kg}$, $[\kappa]=\mathrm{s^{-2}}$.\\
\bottomrule
\end{longtable}
A normálirányhoz kapcsolódó mennyiségek előjele a $H$ választásától függ. A grazing-fejezetben a megengedett tartomány $H>0$; a Filippov-fejezetben $F_1$ a $H>0$, $F_2$ a $H<0$ oldali mező.


\section*{Egységes szakkifejezések}
A speciális fogalmakra az alábbi angol elnevezéseket használjuk. A bevett magyar matematikai és mechanikai szavak (például egyensúly, kapcsolófelület, sajátérték, tapadás, csúszás) megmaradnak. A \emph{sliding} az állapot kapcsolófelületen végzett Filippov-mozgása; a mechanikai csúszás ettől külön fogalom. A szalagon mozgó test mechanikai tapadása például sliding mozgásnak felelhet meg az állapottérben.
\begin{longtable}{@{}p{.39\textwidth}p{.55\textwidth}@{}}
\toprule Egységes elnevezés & Jelentés a jegyzetben\\\midrule\endhead
piecewise-smooth & Tartományonként sima rendszer; a kapcsolófelületen külön illesztési szabály lehet szükséges.\\
sliding; sliding region & A kapcsolófelületen végzett mozgás, illetve ennek tartománya.\\
attracting / repelling sliding & A két oldali mező a felület felé, illetve attól elfelé mutat.\\
sliding mode control & A kapcsolófelület elérésére és az azon előírt dinamikára épülő szabályozás.\\
stick--slip & Mechanikai tapadási és csúszási szakaszok váltakozása.\\
grazing & Pálya és kapcsolófelület érintkezése nulla normálsebességgel.\\
square-root mapping & Négyzetgyökös korrekciót tartalmazó leképezés.\\
reset leképezés & Az esemény előtti állapotot az esemény utáni állapotba vivő $R$ leképezés.\\
saltation-mátrix & Azonos referenciaidőre hozott perturbációk eseményen átvezető lineáris leképezése.\\
discontinuity-induced bifurcation & A kapcsolási határ eléréséhez kapcsolódó bifurkációk gyűjtőneve.\\
border-collision bifurcation & Diszkrét leképezés fixpontjának találkozása az ágakat elválasztó határral.\\
boundary equilibrium bifurcation & Folytonos idejű rendszer egyensúlyának találkozása a kapcsolófelülettel.\\
admissible / virtual equilibrium & A saját mező érvényességi tartományába eső / azon kívül fekvő egyensúlyjelölt.\\
pseudo-equilibrium & A sliding dinamika egyensúlya.\\
saddle-node / non-smooth fold & Két egyensúlyi ág találkozása sima, illetve a tárgyalt nemsima esetben; a két mechanizmust megkülönböztetjük.\\
period-doubling bifurcation & Kétszeres periódusú ciklus megjelenésével járó bifurkáció.\\
pseudo-arclength continuation & Megoldáság követése kibővített rendszerrel és érintőhöz kötött kiegészítő feltétellel.\\
shooting / multiple shooting & Kezdetiérték-feladatok integrálására épülő peremfeladat-megoldás egy, illetve több szegmenssel.\\
\bottomrule
\end{longtable}
